% sec:schw-kruskal — Kruskal–Szekeres and Penrose Diagrams
\section{Kruskal--Szekeres and Penrose Diagrams}
\label{sec:schw-kruskal}

Neither Schwarzschild coordinates nor Eddington--Finkelstein coordinates
cover the full geometry.  The ingoing patch sees the exterior and the
black hole interior but misses the white hole; the outgoing patch sees
the exterior and the white hole but misses the black hole interior.
A natural question arises: is there more spacetime beyond what these
coordinate systems reveal?  The answer matters even for the ray tracer,
which operates only in the exterior --- understanding the global
structure is what confirms that a photon crossing the horizon can never
return.  To see the complete picture, we need coordinates that cover the
\emph{maximal analytic extension} of the Schwarzschild spacetime.

\paragraph{Kruskal--Szekeres coordinates.}

The strategy is to combine the ingoing and outgoing null coordinates
into a pair that is everywhere nonsingular.  Recall the advanced and
retarded times from \cref{sec:schw-coordinates}:
$v = t + r_*$ and $u = t - r_*$, where $r_*$ is the tortoise coordinate
(\cref{eq:tortoise}).  In terms of $u$ and $v$, the Schwarzschild line
element outside the horizon becomes
%
\begin{equation}
  \d s^2 = -\!\left(1 - \frac{2M}{r}\right)\d u\,\d v
           + r^2\,\d\Omega^2 \,,
  \label{eq:ks-double-null}
\end{equation}
%
where $r$ is determined implicitly by $v - u = 2r_*$.  The prefactor
$(1 - 2M/r)$ still vanishes at the horizon.  To remove it, define
%
\begin{equation}
  U = -e^{-u/4M} \,, \qquad V = e^{v/4M} \,.
  \label{eq:ks-UV-def}
\end{equation}
%
These are simply exponential rescalings of the null coordinates,
chosen so that the $1/4M$ in the exponent cancels the $4M$ that
appears when differentiating the tortoise coordinate.  Differentiating
gives $\d U = (4M)^{-1}\,e^{-u/4M}\,\d u$ and
$\d V = (4M)^{-1}\,e^{v/4M}\,\d v$, so
$\d U\,\d V = (4M)^{-2}\,e^{(v-u)/4M}\,\d u\,\d v$.
The product of the exponentials is
%
\begin{equation}
  UV = -e^{(v-u)/4M} = -e^{r_*/2M}
     = -e^{r/2M}\!\left(\frac{r}{2M} - 1\right) ,
  \label{eq:ks-UV-product}
\end{equation}
%
where the last step substitutes the definition of $r_*$
(\cref{eq:tortoise}) and is valid for $r > 2M$.  The right-hand
side of \cref{eq:ks-UV-product} is smooth for all $r > 0$, so we
adopt it as the implicit definition of $r(U,V)$ on the full
extension.  The key identity is that the zero of $(1 - 2M/r)$
at the horizon is matched by a zero of $UV$: the exponential
rescaling converts a singular coordinate factor into a smooth
conformal prefactor.

Rearranging, $e^{(v-u)/4M} = -UV$ and
$(1 - 2M/r) = -(2M/r)\,UV\,e^{-r/2M}$.  Substituting into
\cref{eq:ks-double-null} and collecting factors yields the
Kruskal--Szekeres line element:
%
\begin{equation}
  \d s^2 = -\frac{32M^3}{r}\,e^{-r/2M}\,\d U\,\d V
           + r^2\,\d\Omega^2 \,,
  \label{eq:ks-metric-null}
\end{equation}
%
with $r(U,V)$ determined implicitly by \cref{eq:ks-UV-product}.
The prefactor $32M^3\,e^{-r/2M}/r$ is strictly positive and finite
for all $r > 0$: the coordinate singularity at $r = 2M$ has been
removed entirely.  The only singularity left is the physical one at
$r = 0$.
\histnote{Kruskal and Szekeres independently discovered these
coordinates in 1960 \cite{Kruskal1960,Szekeres1960};
Fuller and Wheeler first analysed the maximal extension in
detail \cite{FullerWheeler1962}.}

For many purposes it is convenient to work with timelike and
spacelike coordinates rather than the null pair $(U, V)$.
Define
%
\begin{equation}
  T = \tfrac{1}{2}(V + U) \,, \qquad
  X = \tfrac{1}{2}(V - U) \,.
  \label{eq:ks-TX-def}
\end{equation}
%
The metric becomes
%
\begin{equation}
  \d s^2 = \frac{32M^3}{r}\,e^{-r/2M}
    \bigl({-}\d T^2 + \d X^2\bigr)
    + r^2\,\d\Omega^2 \,,
  \label{eq:ks-metric}
\end{equation}
%
where $r$ is now defined implicitly by
%
\begin{equation}
  X^2 - T^2 = \left(\frac{r}{2M} - 1\right)e^{r/2M} \,.
  \label{eq:ks-r-implicit}
\end{equation}
%
The conformal factor multiplying $(-\d T^2 + \d X^2)$ is positive
for $r > 0$, so light cones in the $(T, X)$ plane make $45^\circ$
angles everywhere --- precisely the feature that makes
Kruskal--Szekeres coordinates valuable.  Causal structure is read
off as easily as in Minkowski spacetime.
\physnote{The $45^\circ$ light cones are not an accident.  The $(T, X)$
block of the metric is conformally flat: a position-dependent rescaling
of the two-dimensional Minkowski metric.  Conformal rescalings preserve
causal structure.}

\paragraph{The four regions.}

Surfaces of constant $r$ satisfy $X^2 - T^2 = \text{const}$:
they are hyperbolas in the $(T, X)$ plane.  The horizon $r = 2M$
corresponds to $X^2 - T^2 = 0$, i.e.\ the pair of lines
$T = \pm X$.  These two null lines divide the $(T, X)$ plane into
four regions:
%
\begin{center}
\begin{tabular}{cll}
  \toprule
  Region & Definition & Physical interpretation \\
  \midrule
  I   & $X > \abs{T}$ & Exterior (our universe) \\
  II  & $T > \abs{X}$ & Black hole interior (future singularity) \\
  III & $X < -\abs{T}$ & Second exterior (disconnected) \\
  IV  & $T < -\abs{X}$ & White hole interior (past singularity) \\
  \bottomrule
\end{tabular}
\end{center}
%
Region~I is the familiar exterior covered by Schwarzschild coordinates
with $t \in (-\infty, \infty)$.  Region~II is the black hole interior:
every future-directed causal curve in this region reaches the
singularity at $r = 0$.  In units where $2M = 1$, the singularity
traces the hyperbolas $T = \pm\sqrt{X^2 + 1}$, bounding Regions~II
and~IV respectively.  Region~III is a second asymptotically flat
exterior,
disconnected from Region~I --- no causal curve connects them.
Region~IV is the time-reverse of Region~II: a white hole from which
matter and light can emerge but never enter.
\warnnote{Region~III is not the ``other side'' of the black hole in
any traversable sense.  A traveller who crosses the future horizon from
Region~I enters Region~II, not Region~III, and inevitably reaches the
singularity.}

Ingoing Eddington--Finkelstein coordinates
(\cref{sec:schw-coordinates}) cover Regions~I and~II but not III
or~IV.  Outgoing Eddington--Finkelstein coordinates
(\cref{eq:ef-outgoing}) cover Regions~I and~IV.
Kruskal--Szekeres coordinates are the only system among those
we've constructed that covers all four regions simultaneously.

\paragraph{Penrose diagrams.}

Kruskal--Szekeres coordinates have a practical limitation: regions at
$r \to \infty$ are pushed to $X \to \infty$, so the diagram extends
infinitely in both directions.  A Penrose (or Carter--Penrose) diagram
solves this by applying a further conformal compactification that maps
the entire spacetime onto a finite region, while preserving the
$45^\circ$ light cones that make causal structure visible.
\histnote{Penrose introduced the conformal compactification technique
in 1963 \cite{Penrose1963}; Carter independently developed the
diagrammatic method for the Kerr family
\cite{Carter1966}.}

The construction begins with the null coordinates $(U, V)$.  Apply
the compactification
%
\begin{equation}
  \tilde{U} = \arctan U \,, \qquad
  \tilde{V} = \arctan V \,,
  \label{eq:penrose-compactify}
\end{equation}
%
which maps $U, V \in (-\infty, \infty)$ to
$\tilde{U}, \tilde{V} \in (-\pi/2, \pi/2)$.  The physical spacetime
does not fill the full square
$(\tilde{U}, \tilde{V}) \in (-\pi/2, \pi/2)^2$: the singularity at
$r = 0$ corresponds to $UV = 1$, which defines the upper and lower
boundaries of the diagram.  In the compactified coordinates, the metric
takes the form $\d s^2 = \Omega^{-2}\,\d\tilde{s}^2$ for a smooth
conformal factor $\Omega$ that vanishes at infinity.  The explicit form
of $\Omega$ involves the Kruskal prefactor and the derivatives of the
$\arctan$ maps; what matters for diagram construction is only that
$\Omega > 0$ everywhere except at conformal infinity, where
$\Omega \to 0$ defines the boundary.  The ``unphysical'' metric
$\d\tilde{s}^2$ is the one actually drawn; since it differs from the
physical metric by a strictly positive function, the causal structure
--- which null curves are future-directed, which events can communicate
--- is identical.

%
\begin{figure}[htbp]
  \centering
  % penrose-schwarzschild.tex — Penrose diagram for the maximally extended
% Schwarzschild spacetime
% Inputable TikZ fragment for fig:penrose-schwarzschild
% Requires: tikz with calc, arrows.meta, decorations.markings (loaded by ehbook.cls)
%
% Shape: diamond with top/bottom vertices flattened into horizontal
% singularity lines.  All null lines at 45°.
%
% Geometry (cm):
%   Half-height h = 1.8 (centre to singularity).
%   Singularity half-width s = 1.0.
%   Spatial-infinity offset w = s + h = 2.8 (ensures 45° edges).
%
\begin{tikzpicture}[
    >={Stealth[length=5pt]},
    font=\footnotesize,
    zigzag/.style={
      decorate,
      decoration={zigzag, segment length=3.5pt, amplitude=1.5pt},
    },
  ]

  % ── Parameters ──────────────────────────────────────────────────────
  \def\h{1.8}    % half-height (centre to singularity)
  \def\s{1.0}    % singularity half-width
  \def\w{2.8}    % spatial infinity offset (= s + h for 45° edges)

  % ── Key points ──────────────────────────────────────────────────────
  \coordinate (ipI)   at ( \s,  \h);   % i+  Region I
  \coordinate (imI)   at ( \s, -\h);   % i-  Region I
  \coordinate (i0I)   at ( \w,    0);  % i^0 Region I
  \coordinate (ipIII) at (-\s,  \h);   % i+  Region III
  \coordinate (imIII) at (-\s, -\h);   % i-  Region III
  \coordinate (i0III) at (-\w,    0);  % i^0 Region III

  % ── Region fills (very faint) ──────────────────────────────────────
  % Region II (top triangle)
  \fill[EHAccretionGold, opacity=0.04]
    (ipIII) -- (ipI) -- (0, 0) -- cycle;
  % Region IV (bottom triangle)
  \fill[EHAccretionGold, opacity=0.04]
    (imIII) -- (imI) -- (0, 0) -- cycle;

  % ── Singularities (zigzag, top and bottom) ─────────────────────────
  \draw[EHHorizonCrimson, very thick, zigzag]
    (ipIII) -- (ipI);
  \draw[EHHorizonCrimson, very thick, zigzag]
    (imIII) -- (imI);

  % Singularity labels
  \node[above, font=\scriptsize, text=EHHorizonCrimson]
    at (0, \h) {$r = 0$};
  \node[below, font=\scriptsize, text=EHHorizonCrimson]
    at (0, -\h) {$r = 0$};

  % ── Null infinity (diamond boundary edges) ─────────────────────────
  % Region I: right side
  \draw[EHMarginGray, thick] (ipI) -- (i0I)
    node[pos=0.5, above right, inner sep=1pt, font=\scriptsize,
         text=EHMarginGray] {$\mathscr{I}^+$};
  \draw[EHMarginGray, thick] (i0I) -- (imI)
    node[pos=0.5, below right, inner sep=1pt, font=\scriptsize,
         text=EHMarginGray] {$\mathscr{I}^-$};
  % Region III: left side
  \draw[EHMarginGray, thick] (ipIII) -- (i0III)
    node[pos=0.5, above left, inner sep=1pt, font=\scriptsize,
         text=EHMarginGray] {$\mathscr{I}^+$};
  \draw[EHMarginGray, thick] (i0III) -- (imIII)
    node[pos=0.5, below left, inner sep=1pt, font=\scriptsize,
         text=EHMarginGray] {$\mathscr{I}^-$};

  % ── Horizons (dashed X through centre) ─────────────────────────────
  \draw[EHHorizonCrimson, dashed, thick]
    (imIII) -- (ipI);  % T = X line
  \draw[EHHorizonCrimson, dashed, thick]
    (ipIII) -- (imI);  % T = -X line

  % ── Region labels ──────────────────────────────────────────────────
  \node[font=\small] at ( 1.4,  0) {I};
  \node[font=\small] at (   0, 0.7) {II};
  \node[font=\small] at (-1.4,  0) {III};
  \node[font=\small] at (   0, -0.7) {IV};

  % ── Conformal boundary labels ──────────────────────────────────────
  % Spatial infinity
  \node[right, font=\scriptsize, text=EHMarginGray]
    at ({\w + 0.08}, 0) {$i^0$};
  \node[left, font=\scriptsize, text=EHMarginGray]
    at ({-\w - 0.08}, 0) {$i^0$};

  % Timelike infinity
  \node[right, font=\scriptsize, text=EHMarginGray, inner sep=1pt]
    at ({\s + 0.08}, \h) {$i^+$};
  \node[right, font=\scriptsize, text=EHMarginGray, inner sep=1pt]
    at ({\s + 0.08}, -\h) {$i^-$};
  \node[left, font=\scriptsize, text=EHMarginGray, inner sep=1pt]
    at ({-\s - 0.08}, \h) {$i^+$};
  \node[left, font=\scriptsize, text=EHMarginGray, inner sep=1pt]
    at ({-\s - 0.08}, -\h) {$i^-$};

  % ── Vertex dots ────────────────────────────────────────────────────
  \fill[EHMarginGray] (i0I)   circle (1.5pt);
  \fill[EHMarginGray] (i0III) circle (1.5pt);
  \fill[EHMarginGray] (ipI)   circle (1.2pt);
  \fill[EHMarginGray] (imI)   circle (1.2pt);
  \fill[EHMarginGray] (ipIII) circle (1.2pt);
  \fill[EHMarginGray] (imIII) circle (1.2pt);

  % ── Sample timelike worldline ──────────────────────────────────────
  % Starts in Region I, crosses future horizon, hits singularity.
  % Must have slope > 1 (steeper than 45°) to be timelike.
  \draw[EHJetBlue, thick, ->]
    plot[smooth, tension=0.6] coordinates {
      (1.6, -0.8)
      (1.2, -0.1)
      (0.7, 0.6)
      (0.3, 1.2)
      (0.15, 1.6)
      (0.1, 1.78)
    };
  \node[right, font=\scriptsize, text=EHJetBlue, inner sep=1pt]
    at (1.6, -0.8) {observer};

\end{tikzpicture}

  \caption[Penrose diagram for the maximally extended Schwarzschild spacetime]{%
    Penrose diagram for the maximally extended Schwarzschild spacetime.
    The entire spacetime fits in a finite diamond; light rays travel at
    $45^\circ$.  The future singularity (top, zigzag) bounds Region~II
    (the black hole interior); the past singularity (bottom) bounds
    Region~IV (the white hole).  The horizons (dashed diagonals) separate
    the four regions.  Spatial infinity $i^0$, future/past timelike
    infinity $i^\pm$, and future/past null infinity $\mathscr{I}^\pm$ are
    marked at the boundary of each exterior region.  A timelike observer
    (solid curve) in Region~I who crosses the future horizon inevitably
    reaches the singularity; no causal curve connects Region~I to
    Region~III.}
  \label{fig:penrose-schwarzschild}
\end{figure}

The resulting diagram has a characteristic diamond shape.  Each of the
four Kruskal--Szekeres regions occupies a triangular quadrant.  The two
singularities at $r = 0$ are horizontal lines (conventionally drawn as
zigzags) at the top and bottom of the diamond.  The horizons are the
$45^\circ$ diagonals connecting the left and right vertices.  The
boundary of the diamond represents conformal infinity.  Decomposing
infinity into distinct pieces lets us state precisely where different
classes of curves end up: massive particles reach timelike infinity,
photons reach null infinity, and spacelike slices terminate at spatial
infinity.  The five standard boundary components for each exterior
region are:
%
\begin{itemize}
  \item $i^0$ (spatial infinity): the right vertex, where $r \to \infty$
    at fixed $t$.
  \item $i^+$ and $i^-$ (future and past timelike infinity): the
    endpoints of future/past-inextendible timelike curves in the
    exterior region.
  \item $\mathscr{I}^+$ and $\mathscr{I}^-$ (future and past null
    infinity): where outgoing and ingoing light rays reach as
    $r \to \infty$.
\end{itemize}
%
\warnnote{The Penrose diagram is a $(1+1)$-dimensional projection:
each point represents a two-sphere of area $4\pi r^2$.  The diagram
faithfully represents radial causal structure but suppresses angular
information.}

\paragraph{Reading causal structure from the diagram.}

The real power of the Penrose diagram is that causal relationships
become geometric.  Two events can be connected by a causal (timelike or
null) curve if and only if one lies inside the light cone of the other
in the diagram.  Three consequences follow immediately:

First, every future-directed causal curve from Region~II terminates at
the future singularity.  The singularity is not a point in space that
one can avoid by steering --- it lies in the future of every event in
Region~II, like a moment in time that everything approaches.
\physnote{Inside the black hole, $r$ becomes a timelike coordinate and
$t$ a spacelike one.  The singularity at $r = 0$ is not a place but a
\emph{time}: the future boundary of the black hole interior.}

Second, no causal curve connects Region~I to Region~III.  To reach
Region~III from Region~I, a future-directed causal curve would have to
cross the future horizon into Region~II.  But once in Region~II, every
future-directed causal curve terminates at the singularity $r = 0$ ---
no path leads onward to Region~III.

Third, the white hole Region~IV is the time-reverse of the black hole
Region~II: matter can leave but not enter.  White holes are
mathematically consistent but widely regarded as astrophysically
irrelevant --- they are unstable under small perturbations and do not
form in gravitational collapse.

\paragraph{Relevance for the ray tracer.}

The ray tracer operates in Region~I (and partially in Region~II for
photons that cross the horizon before termination).  The maximal
extension is therefore not directly needed for rendering.  Its value
is conceptual: before this construction, it might seem that a photon
hitting $r = 2M$ encounters a barrier.  The Kruskal--Szekeres and
Penrose diagrams show the opposite: the horizon is an interior surface,
and the true boundary is the spacelike singularity at $r = 0$.  This
is why the codebase's termination condition tests $r < r_{\min}$ rather
than $r \leq 2M$ --- crossing the horizon is unremarkable; approaching
the singularity is where the physics (and the numerics) break down, as
curvature invariants diverge and the integration loses accuracy.
\xrefnote{The termination conditions and their numerical
implementation are discussed in \cref{sec:schw-orbits}.}

For the Kerr black hole, the causal structure is considerably richer:
the maximal extension has an infinite chain of asymptotic regions
connected by inner horizons, and the singularity is a ring rather
than a point.  The Penrose diagram technique generalises directly,
and we return to it in \cref{sec:kerr-horizons}.
